\chapter{Os Preceitos}

\lettrine{E}{sta noite é lua nova}\footnote{Um dia de vigilância lunar.} e hoje reafirmamos o
nosso compromisso para com \emph{sila}: para os bhikkhus é o \emph{Pāṭimokkha}, para as
monjas os dez preceitos, e para os Anagārikas os oito preceitos. Ao reafirmar o
nosso compromisso, podemos observar estes oito preceitos num nível sublimado de
interpretação. Assim, com o primeiro preceito -- \emph{pāṇātipātā veramaṇī} --
abstermo"-nos de matar outras criaturas -- mesmo que nenhum de nós seja propenso
a assassinar ou à violência física, é importante deixar claro na nossa mente,
que a nossa intenção nesta vida é inclusivamente não magoar os outros
intencionalmente. Em seguida, o segundo preceito -- \emph{adinnādānā veramaṇī}
-- enceta não apenas abstermo"-nos de roubar, mas respeitar a propriedade dos
outros; não perturbar ou fazer mau uso do que pertence aos outros. Esta é uma
maneira de tornar isto bem definido em nossa consciência.

\emph{Abrahmacariyā veramaṇī}, o terceiro preceito, é o voto do celibato.
Estamos num momento em que há tanta preocupação com a SIDA e as doenças
venéreas. Um abuso total da sexualidade ocorreu nas últimas décadas, quando as
pessoas foram totalmente irresponsáveis e buscaram prazer nas actividades
sexuais sem ter em conta as consequências. O resultado é que agora temos dilemas
morais sobre o aborto e as várias doenças e os problemas que surgem e como
resolvê"-los. O que devíamos fazer? Tentar promover o uso de preservativos e
outras medidas profilácticas, para que as pessoas possam fazer tudo o que
quiserem sem terem de se conter? Ou promover pílulas e dispositivos para
prevenir a gravidez para que ninguém tenha que escolher entre fazer um aborto ou
ter um bebé? Em suma, o que nunca é mencionado é qualquer tipo de posição moral.
Parece ser algo que simplesmente não se menciona. O celibato nunca é sequer
considerado como um modo de vida possível.

Mas, realmente, quando consideramos a nossa vida como seres humanos, há uma
maneira de viver mais inteligente. Podemos assumir a responsabilidade pela nossa
existência e abstermo"-nos de envolver outras pessoas ou mesmo de explorar o
nosso próprio corpo para a busca desse tipo de prazer. Seguir o preceito do
celibato é bastante enobrecedor. Eleva"-nos: ser celibatário é um potencial, uma
possibilidade para desenvolver a meditação pelo autodomínio necessário à
realização da verdade. O celibato é algo que tem que ser a própria pessoa a
assumir, não é algo que possa ser forçado; isso já não seria castidade; seria
tirania. Tem que ser uma opção nossa, algo a que nos elevamos como seres
individuais. Não queremos regressar a uma posição puritana de ``Não farás'',
ameaçando pessoas com 84.000 éones em infernos de fogo, ardendo em dor absoluta
por causa de qualquer tipo de prazer sexual. Não estamos a tentar trazer medo à
mente ou a intimidar, mas sim a encorajar o que há de nobre e belo na nossa
humanidade.

Suponho que sejam capazes de se motivarem e, por isso, apresento esta
oportunidade para praticar. Às vezes, as pessoas podem ter opiniões muito baixas
sobre si mesmas, o que não é realmente verdade. Talvez nunca tenham tido uma
oportunidade ou nunca sentiram que alguém confiava nelas o suficiente para se
motivarem. Estamos a tentar trazer para a nossa vida monástica esse tipo de
valor, esse tipo de beleza, de modo a que o monasticismo seja algo que é ``belo
no início, belo no meio e belo no fim'' e não um tipo de tirania imposta ou uma
marcha forçada.

Precisamos assumir essa responsabilidade por nós mesmos, em vez de entregá"-la a
outra pessoa, à espera de que outra pessoa nos ilumine, ame, conduza ou
repreenda. O potencial espiritual de cada ser, deve ser aqui reconhecido. Temos
essa capacidade maravilhosa de nos elevar"-nos às coisas em vez de nos
afundarmos.

``Elevar'' não é uma força intencional; é a capacidade de ir além da inércia ou
das tendências habituais da vida em direcção a algo mais elevado; é aprender a
simplesmente prestar atenção à respiração ou ser mais paciente, mais compassivo,
mais amável -- connosco próprios e com os outros. Tudo isto é o esforço de se
elevar e agir no momento certo. Isto não significa que temos que ter sempre
sucesso ou provarmos a nós mesmos; significa elevarmo"-nos para enfrentar uma
situação de forma inteligente, com consciência e sabedoria. E esta é uma
possibilidade para nós: não temos que ser apanhados pela força do hábito e
perdidos no reino da ilusão.

Com respeito à fala, \emph{musāvādā veramaṇī}, o preceito é abstermo"-nos de
falar incorrectamente: e é tão fácil cairmos na visão de nós próprios se usarmos
as formas vulgares de falar ``eu sou'' ou ``pobre de mim''! Observemos a forma
como o Buddha usou a linguagem: ``há'' sofrimento, ``há'' raiva / ganância /
ilusão. Isto é um exemplo de como nos refrearmos de discurso incorrecto. Se
começarmos a reflectir dessa forma, afectará como vemos as coisas. Nesta
comunidade, temos uma vontade de aprender a comunicar e procuramos falar de uma
forma clara e honesta, mas não reivindicativa ou iludida. Em contraste, na
sociedade, uma pessoa tenta ser inteligente no discurso, arguta, divertida -- e,
com uma mente inteligente, os hábitos de falar da pessoa podem ser bastante
cruéis e maus. Mas desistimos disso e tentamos usar a fala como algo belo, claro
e sem criar visões erradas aos outros.

\emph{Musāvādā veramaṇī} não é só abster"-se de mentir, mas envolve a intenção de
assumir a responsabilidade de falar. Toda esta função da nossa humanidade tem
muito de milagre quando a contemplamos. E, no entanto, simplesmente a assumimos
como garantida. Podemos usar o nosso discurso para dizer piadas nojentas,
palavrões, para praguejar e coscuvilhar, insultar e todos os tipos de coisas
mesquinhas, horríveis e desonestas. Ou, podemos respeitar esse dom maravilhoso
que temos e aprender a usá"-lo de uma maneira bela, precisa e amável.

E então com \emph{Surāmeraya"-majja"-pamādaṭṭhānā veramaṇī} -- abstendo"-se de
intoxicantes: reflicta"-se em quão afortunados somos por não termos que beber,
usar drogas e injectar heroína. Isso afecta todos os níveis da sociedade.
Homens, mulheres e crianças -- todas as raças, todas as classes estão a ser
apanhadas nas garras dessas drogas viciantes. Também existem os cigarros e
álcool -- tudo prejudicial e ilusório para a mente humana. Quando ficamos
enublados com drogas e bebidas, não podemos ser responsáveis pelo que dizemos,
podemos?

Lembro"-me de quando bebia, era para não ter que ser responsável pelo que dizia!
Quando intoxicados, perdemos o senso de timidez e vergonha em relação à conduta
sexual. Bebe"-se umas quantas bebidas e, de repente, muitas inibições
simplesmente desaparecem. Eu não gostava de matar pessoas, mas certamente não
hesitei em me livrar de insectos irritantes e outras coisas que não gostava.
Podia ver"-se, sob o efeito das drogas e da bebida, o quão facilmente o sentido
de decência moral e compromisso podem desaparecer. Hoje em dia encontramos
jovens a prostituir"-se para conseguir dinheiro para comprar bebida e drogas --
até mesmo pessoas com doze ou treze anos, aqueles a quem chamávamos de crianças!

Depois, há os preceitos renunciantes, aqueles que simplificam as nossas vidas.
Abster"-se de comer depois do meio"-dia, de se divertir e adornar"-se. Para os
seres humanos, existe todo um reino de diversão e entretenimento disponível
através de comer, dançar, cantar, jogos, filmes, TV e espectáculos. Depois temos
o sono. Há a tentação de se passar imenso tempo à procura de conforto e dormir.
Ora, estas coisas não são imorais, são? Eu não estou a dizer que jantar seja uma
actividade imoral -- ou dançar e cantar, convenhamos -- mas estamos a tentar
conter"-nos e abster"-nos de oportunidades de nos distrairmos com os prazeres
sensoriais, de modo a que possamos observar e reflectir.

Estas são normas e preceitos para reflexão e não regras de Deus. Não devem ser
observadas da posição ``Não farás'' (Bíblia). Cada um dos preceitos é uma
resolução, algo que estamos a assumir e não algo que Deus nos está a impor.
Então elevamo"-nos a esses preceitos e tomamos uma resolução de forma a
possuirmos isso nas nossas mentes quando somos tentados a agir sob os impulsos
que podemos estar a experienciar. Afinal, a maioria de nós vimos de passados
bastante indulgentes, onde nunca fomos realmente encorajados a dominarmo"-nos.

\emph{Sīla} é uma oportunidade digna e maravilhosa que temos como seres humanos.
Escolhemos ser morais. Não estamos a ser morais porque temos medo de ser
imorais. Escolhemos fazer isto e elevamo"-nos ao que é nobre, bom, amável e
generoso.

Reconhecidamente, as atracções mundanas continuam fortes e não faz parte do meu
ensino condená"-las. Não sou contra a vida mundana, nem estou a tentar elevar o
monasticismo como algo que todos deveriam fazer. Também conseguimos viver uma
vida mundana muito boa e saudável -- ser saudável não é só uma prerrogativa
exclusiva de monges e monjas, certo? Às vezes os leigos pensam que sou um
``monge fanático'' porque enfatizo o valor deste modo de vida.

Mas a atitude aqui é de reflexão, ao invés de remoer as coisas com
ressentimentos ou de adoptar uma posição acerca de qualquer coisa. O nosso
objectivo é de desenvolver esta capacidade reflexiva da mente -- e as convenções
específicas que usamos são desenvolvidas em torno disso. É disto que se trata o
ensinamento do Buddha. Toda a convenção do Vinaya (disciplina) e dos
ensinamentos do Dhamma existe para nos ajudar nesse sentido. Algumas
pessoas dizem que é possível fazer isso como leigo, e isso não se pode negar;
mas se não se consegue usar o estilo de vida que é deliberadamente estabelecido
para o Dhamma"-Vinaya, o que fará alguém pensar que consegue fazê"-lo
noutra convenção que seja? É isso que eu vos quero levar a ver. Olhem para vós
próprios; nesse querer algo que não se tem ou querer fugir do que se tem.
Simplesmente observem isso dentro de vós, essa inquietação, descontentamento e
movimento da vossa mente.

Às vezes, é claro, uma pessoa ainda não quer desistir, ainda quer renascimento e
felicidade e coisas mundanas. É justo! Mas eu não quero que andem por aí a
mentir a vós mesmos. Se quiserem seguir o vosso próprio caminho, ter
renascimentos e felicidade mundana, então essa é a vossa decisão -- mas não se
iludam pensando que estão a fazer outra coisa. Porque, se realmente entendem os
ensinamentos do Buddha, então não há para onde ir e nada para fazer. É assim que
as coisas são. Estamos a viver uma vida que é para esse tipo de reflexão. Então
podemos observar esse desejo de estar noutro lado como um movimento da mente,
vê"-lo e reconhecê"-lo pelo que é. Quer o façamos ou não, depende da vontade de
cada um!

Permitam"-se morrer para o momento. Investiguem e observem como as coisas são.
Tudo o que surge cessa. No todo, tudo se encaixa nesse padrão, não é? Desta
forma, podemos apenas reflectir sobre a banalidade mundana do dia-a-dia das
nossas vidas. Já que não podemos dançar e cantar, ir a espectáculos, bares,
jogos de futebol, restaurantes e seguir os prazeres e distracções do mundo,
então a normalidade das nossas vidas assume mais significado. Se estamos
habituados a um elevado nível de excitação, as coisas comuns são simplesmente
chatas e estaremos sempre em busca de alguma nova emoção ou experiência. O
monasticismo é um modo de vida entediante, apenas uma rotina. Não temos como
objectivo ter coisas emocionantes para fazer, ou distracções, porque na
meditação estamos cientes de como as coisas normalmente são na consciência.
Portanto, já não estamos no encalce nem a tentar encontrar, o extraordinário. É
somente através do apaziguamento e do domínio, não por seguirmos a inquietação e
sermos tomados pelas emoções, que temos a oportunidade de realizarmos o
Incondicionado. Só quando conseguimos abrir mão, acalmarmo"-nos e reflectirmos, é
que há uma compreensão do fim do reino condicionado -- em que tudo o que surge
cessa -- e uma compreensão do Nibbāna. Não há forma de realizarmos o
Nibbāna pelo esforço e tentativa de se atingir e alcançar, e sermos
apanhados no lado do surgir. Temos que abandonar isso.

A concretização do abandonar"-se aquilo que surge na mente, conduz ao testemunhar
da cessação do que surgiu. Então, há a verdadeira paz de permitir que as coisas
sejam como são. Já não somos alguém que tem que chegar a algum lugar, fazer
algo, livrar"-se de algo ou mudar algo. Quando somos apanhados a distrair"-nos com
prazeres, logo somos alguém, e alguém que tem que encontrar a felicidade ou ter
sucesso ou se tornar"-se algo. Não importa quanta excitação e prazeres possamos
experimentar, temos que obter mais do que isso. Nunca estamos satisfeitos com a
excitação e as aventuras da vida. Estas apenas nos fazem sermos apanhados
naquele movimento de ter que se ter mais e mais -- até ficarmos exaustos. Então
vamos para o extremo oposto, onde, por estarmos cansados e exaustos de toda a
excitação e estimulação, simplesmente desabamos, caímos no sono, nos drogamos ou
embriagamos. Não queremos existir. Só suportamos determinada excitação e depois
não conseguimos aguentar mais.

Estar"-se excitado continuamente é uma dimensão infernal, não é? Na casa da minha
irmã na Califórnia, eles têm todos aqueles canais de televisão e relés
transmissores de cabo. Podemos sentar"-nos e ter cerca de setenta canais à nossa
disposição o tempo todo. As pessoas ganham o hábito de apenas mudarem de canal
se alguma coisa ficar um pouco entediante, um pouco desinteressante ou
desagradável. Mudam simplesmente para o próximo -- há sempre um tiroteio ou uma
linha de refrão para excitar uma pessoa. É uma espécie de reino infernal -- é
desagradável ter uma mente que precisa de ser estimulada um momento após o
outro.

Então vemos a beleza de uma vida que é baseada na serenidade, domínio moral,
nobreza, generosidade, bondade e reflexão no Dhamma. É maravilhoso poder
ter a oportunidade e o incentivo de contemplar a própria existência e treinar de
forma a respeitarmo"-nos a nós mesmos. Podemos dirigir"-nos para sermos uma pessoa
contente e alegre, em vez de uma ávida e obcecada.

\photoPage{image-005-cottage-alms-round.jpg}{Circumambulação no Vihara de Ratanagiri}{Foto por Ron Livingston}
